% =====================================================================
%  capitulos/cap08.tex — Capítulo 8. La estructura eventiva
%  Autora: Mabel Giammatteo
%  Fuente: Capitulo_08_normalizado.docx
% =====================================================================
\capitulo{La estructura eventiva}{Mabel Giammatteo}

\section{Introducción}

Según hemos visto en el \capref{6}, la existencia misma de la oración depende de
su anclaje temporal respecto del momento de la emisión. Pero, además del tiempo,
la organización de la oración obedece a otra categoría que, desde el punto de
vista gramatical, en español se da casi siempre amalgamada morfológicamente con
el tiempo, pero que, como veremos, tiene una gran incidencia sintáctica: el
aspecto. A diferencia del tiempo, que es más externo, el aspecto tiene que ver
con la manera en que un evento se desarrolla u ocurre ---con o sin cambio; con o
sin límite; de forma única o repetida; en un período, acotado o no, o en un
instante; etc.---.

En este capítulo veremos cómo las propiedades más relevantes del aspecto
---duración, delimitación, dinamismo, reiteración--- han permitido identificar
los distintos subtipos de eventos que los verbos expresan. Esta clasificación
resulta de gran utilidad porque, como sostienen \textcite{bosque2009}, sirve
para explicar las propiedades combinatorias de los verbos con sus complementos
y adjuntos. Es necesario tener en cuenta, como plantean los mencionados autores,
que «[a]lgunos verbos pertenecen unívocamente a una sola clase, pero la mayor
parte puede pertenecer a varias si el contexto sintáctico lo permite»
\parencite[307]{bosque2009}. Por eso, también nos ocuparemos de mostrar cómo la
información aportada por el lexema verbal puede ser modificada por la que
aportan sus complementos y modificadores, o incluso por el tiempo y modo verbal
en que se encuentra.

\section{El evento}\label{sec:08-evento}

Informalmente, un \termino{evento} es cualquier clase de situación o suceso
denotado por un predicado \parencite[158]{pustejovsky2019}. A diferencia de los
objetos, que se ubican en el espacio, como \forma{el libro} en
\exref{ex:08-libro}, los eventos, como \forma{ir al cine} en \exref{ex:08-cine},
ocurren y transcurren en el tiempo.

\begin{exe}
  \ex\label{ex:08-libro} El libro está sobre la mesa.
  \ex\label{ex:08-cine} Fuimos al cine ayer a la noche.
\end{exe}

Si bien el tiempo y el aspecto son dos categorías que tradicionalmente se han
relacionado con el evento, ambas manifiestan nociones diferentes:

\begin{destacado}
Mientras el \termino{tiempo} relaciona al evento con una referencia externa que
da cuenta de su anterioridad, simultaneidad o posterioridad con respecto al
momento de la emisión, el \termino{aspecto} constituye una determinación más
interna, que provee información sobre el evento en sí, es decir, sobre cómo
ocurre o transcurre, si se repite o cuánto se extiende.
\end{destacado}

Actualmente se reconocen dos formas de expresar el aspecto:

\begin{itemize}
  \item El \termino{aspecto léxico} o \termino{de situación}, que transmite
        información sobre el tipo de evento de que se trata y se manifiesta en
        la raíz verbal.
  \item El \termino{aspecto gramatical} o \termino{flexivo}, que permite al
        hablante seleccionar el punto de vista desde el que presenta al evento y
        que es manifestado por el verbo en su desinencia.
\end{itemize}

\subsection{Las propiedades de los eventos}\label{sec:08-propiedades}

Si bien ambas formas de expresar la aspectualidad interactúan en la oración,
para la clasificación de los eventos, en principio, partimos del aspecto léxico.
Siguiendo a \textcite{pustejovsky2019}, algunas de las propiedades relevantes de
esta categoría, que inciden en la clasificación de eventos que veremos más
adelante, son las siguientes:

\begin{itemize}
  \item La \termino{duratividad}, que distingue entre eventos instantáneos o de
        duración mínima ---\forma{entrar, disparar} (un arma),
        \forma{encontrar}--- y eventos durativos, cuyo desarrollo implica un
        cierto tiempo ---\forma{estudiar, escribir, comer}---.
  \item La \termino{delimitación}, que diferencia entre eventos que terminan
        ---\forma{Juan construyó una casa en diez meses}--- y eventos que se
        extienden por un período de tiempo y que, en principio, no implican un
        límite ---\forma{Juan ama a su mujer}---.
  \item La \termino{dinamicidad}, que tiene en cuenta si los eventos producen
        cambios o no. Los estados, como \forma{amar, vivir} o \forma{ser alto},
        no involucran transformaciones y no pueden limitarse; así, no podemos
        decir \juicio{*}{\forma{Juan es alto hasta mañana}}. A diferencia de esto, los eventos dinámicos producen
        cambios a través del tiempo: \forma{Juan pintó la pared} (el evento
        produjo un cambio en el estado físico de la pared).
  \item La \termino{telicidad}, según la cual algunos eventos dinámicos tienen
        «una culminación natural» (un \textit{telos}\footnote{La palabra
        \textit{telos} proviene del griego τέλος, \sig{fin}, \sig{objetivo},
        \sig{propósito} o \sig{meta}.} interno), que impide que se prolonguen;
        así, por ejemplo, \forma{comer una manzana} o \forma{fumar un cigarrillo}
        son eventos télicos que no admiten continuarse: \forma{Juan comió una
        manzana (*y siguió haciéndolo toda la tarde)}.\footnote{Por convención,
        en los ejemplos gramaticales se marcan entre paréntesis las secuencias
        optativas. Si se coloca un asterisco antes del paréntesis inicial, la
        secuencia incluida se considera requerida y no puede omitirse. Si, en
        cambio, el asterisco aparece después del paréntesis inicial, se lee que
        agregar la secuencia delimitada por los paréntesis a la oración produce
        agramaticalidad.} En cambio, \forma{Juan leía} es un evento atélico (no
        télico) porque no tiene una culminación en sí mismo y acepta ser
        continuado: \forma{Juan leía y siguió haciéndolo toda la tarde}.
  \item La \termino{iteración}, que implica que algunos eventos son únicos y se
        producen una sola vez ---\forma{Ayer fui al cine} o \forma{Echó un
        tremendo bostezo}--- y otros se repiten ---\forma{Siempre voy al cine los
        fines de semana; Bostezaba sin parar} (\sig{daba varios bostezos})---.
  \item La \termino{intensidad}, que se refiere a que los eventos pueden tener
        diversos grados de fuerza o énfasis en su realización, como, por
        ejemplo, \forma{lloviznar / llover / diluviar} (\sig{llover
        intensamente}) o \forma{picotear} (algo comestible) / \forma{comer /
        devorar}.
\end{itemize}

\subsection{Las clases de eventos}\label{sec:08-clases}

La tradición gramatical adjudica a Aristóteles una primera diferenciación entre
verbos que describen eventos que tienen un punto final, a los que hemos
denominado télicos, y eventos que carecen de término final. En la actualidad, la
clasificación de eventos más reconocida proviene de otro filósofo,
\textcite{vendler1967}, quien distinguió cuatro subclases, a las que luego
\textcite{comrie1976} añadió una quinta. Veamos algunos ejemplos de cada clase
para luego aportar la caracterización de cada una de ellas, a partir de las
propiedades identificadas en §\ref{sec:08-propiedades}:

\begin{itemize}
  \item Estados: \forma{María tiene hambre; Luis es alto}.
  \item Actividades: \forma{Ana juega a las cartas; Corren alrededor de la
        plaza}.
  \item Realizaciones: \forma{Ana cruzó la calle; Comió un pedazo de torta}.
  \item Logros: \forma{La bomba explotó; Encontró lo que buscaba}.
  \item Semelfactivos: \forma{Bostezaba sin parar; Tosía estrepitosamente}.
\end{itemize}

Los \termino{estados} son sucesos durativos [+dur] y, por su carácter estático y
sin cambios [$-$din] ---\forma{ser alto, temer, estar triste}---, se diferencian
de las otras clases, que dan cuenta de situaciones con cambio o dinámicas
[+din], cuyo desarrollo puede implicar fases o momentos diferenciados. Una
característica adicional, relacionada con la falta de dinamismo de estas
predicaciones, es que su mantenimiento no implica ningún esfuerzo deliberado,
por lo que se consideran no agentivas.

Las \termino{actividades} o \termino{procesos} son eventos durativos [+dur] y
dinámicos (implican cambio) [+din] ---\forma{trabajar, correr, leer}---. Al
igual que los estados, su culminación no involucra un fin o \textit{telos}
interno [$-$tél], por lo que se distinguen de las realizaciones [+tél] y los
logros [+tél]. Las actividades no son télicas porque no necesitan llegar a su
culminación para haber tenido lugar o haberse realizado, ya que «se cumplen en
cada una de sus partes» \parencite[62]{garciafernandez1999}. Así, si decimos que
\forma{Ana juega a las cartas} y deja de hacerlo, igualmente puede decirse que
ha jugado a las cartas; en cambio, cuando decimos que \forma{Julia cruza la
calle}, si el evento es interrumpido, no se puede afirmar que cruzó la calle.

Las \termino{realizaciones} son eventos durativos [+dur] y dinámicos [+din], que
tienen una culminación inherente vinculada a la entidad afectada por el evento,
por lo general manifestada por el COD,\footnote{Ampliaremos esta caracterización
en §\ref{sec:08-complementos}.} y, además, se consideran eventos télicos [+tél]
---\forma{leer una carta, escribir un e-mail, construir un puente}---. Como se
mostró con el ejemplo \forma{Julia cruza la calle}, si la realización es
interrumpida, no puede decirse que efectivamente ocurrió.

Los \termino{logros} son eventos dinámicos [+din] y télicos [+tél], como las
realizaciones, pero, como son de cumplimiento instantáneo, se considera que
tienen una duración despreciable o que suceden instantáneamente [$-$dur]
---\forma{llegar, morir, hervir, encontrar}---.

Los \termino{semelfactivos} son eventos instantáneos [$-$dur] y dinámicos
[+din] que, a diferencia de los logros, no tienen \textit{telos} [$-$tél], lo
que permite que sean repetibles. «Usualmente expresan una ocurrencia única de
eventos que suelen ocurrir en secuencias» \parencite[161]{pustejovsky2019}
---\forma{toser, golpear, pestañear}---.

\begin{ejercicio}\label{ej:08-clases}
Lea atentamente el siguiente fragmento y luego resuelva lo pedido:

\begin{consignas}
  \item Según se muestra en los ejemplos resueltos de la tabla, identifique la
        clase eventiva a la que pertenece cada uno de los verbos.
  \item Complete la tabla indicando las características aspectuales de cada
        verbo.
\end{consignas}

\begin{quote}\small
Indiana Jones es un arqueólogo que viaja por el mundo buscando tesoros valiosos
para la humanidad. En \textit{Los cazadores del arca perdida}, que dirigió Steven
Spielberg, Indiana Jones vive una serie de aventuras hasta que finalmente
recupera el arca de la alianza. En sus andanzas, Jones enfrenta rivales a los
que golpea hasta que los derriba o atrapa con su lazo. Sin embargo, este
intrépido personaje tiene una terrible fobia que adquirió en su juventud: las
serpientes. Según aseguran, la saga del popular profesor Jones tiene un nuevo
guionista que ya ha escrito el libro para la próxima película.
\end{quote}
\pendiente{El original dice «Stephen Spilberg»}

\begin{center}\small
\begin{tabular}{@{}lcccl@{}}
  \toprule
  Verbo & Duratividad & Dinamicidad & Telicidad & Clase eventiva \\
  \midrule
  \forma{ser arqueólogo}      & + & $-$ & $-$ & estado \\
  \forma{viajar}              & + & + & $-$ & actividad \\
  \forma{buscar}              & & & & \\
  \forma{dirigir una película} & & & & \\
  \forma{vivir}               & & & & \\
  \forma{recuperar}           & & & & \\
  \forma{enfrentar}           & & & & \\
  \forma{golpear}             & & & & \\
  \forma{derribar}            & & & & \\
  \forma{atrapar}             & & & & \\
  \forma{tener}               & & & & \\
  \forma{adquirir}            & & & & \\
  \forma{asegurar}            & & & & \\
  \forma{escribir el libro}   & & & & \\
  \bottomrule
\end{tabular}
\end{center}
\end{ejercicio}

\section{Factores que intervienen en la delimitación del evento}\label{sec:08-factores}

Un aspecto importante a tener en cuenta en relación con la clasificación de los
eventos es que la información contenida en el lexema verbal puede ser
modificada por la que aportan otros constituyentes de la oración, como el
sujeto, los complementos y modificadores del verbo, así como también el punto de
vista manifestado por el aspecto flexivo o los valores temporales o modales
aportados por la desinencia verbal. Veamos las distintas posibilidades.

\subsection{Los complementos del V}\label{sec:08-complementos}

La presencia del COD resulta esencial para delimitar el evento, como lo muestra
la oposición entre \forma{escribir} (en el sentido de \sig{ser escritor})
\exref{ex:08-gm1} y \forma{escribir una novela} \exref{ex:08-gm2}. El primer
evento, al ser no delimitado, constituye una actividad; en cambio, el segundo,
al estar acotado en su realización al tiempo requerido para la escritura de la
novela, se convierte en una realización.

\begin{exe}
  \ex\label{ex:08-gm1} García Márquez \textit{escribió} durante toda su vida.
  \ex\label{ex:08-gm2} García Márquez \textit{escribió} [la novela \textit{Cien
      años de soledad}].
\end{exe}

Sin embargo, no todo COD delimita el evento: solo lo hacen aquellos que, como
\forma{la novela Cien años de soledad} en \exref{ex:08-gm2}, sirven para medir
el tiempo de desarrollo; en cambio, los objetos que no experimentan un proceso
de cambio o movimiento que los afecte no delimitan al evento. Así,
\exref{ex:08-cuadro} y \exref{ex:08-auto}, en que los objetos \forma{el cuadro} y
\forma{el auto de su padre} no sufren ninguna modificación debido al evento, se
consideran actividades.

\begin{exe}
  \ex\label{ex:08-cuadro} El estudiante \textit{contemplaba} [el cuadro]
      extasiado.
  \ex\label{ex:08-auto} Mariano \textit{manejó} [el auto de su padre] con gran
      precaución.
\end{exe}

Asimismo, tampoco en \exref{ex:08-carne} y \exref{ex:08-relatos} la presencia
del COD es suficiente para delimitar los respectivos predicados. Pero ¿por qué?
¿Qué sucede en estos casos?

\begin{exe}
  \ex\label{ex:08-carne} \textit{Vende} [carne] al exterior.
  \ex\label{ex:08-relatos} García Márquez \textit{escribió} [relatos] a lo largo
      de su vida.
\end{exe}

Los casos como los de \exref{ex:08-carne} y \exref{ex:08-relatos} siguen siendo
actividades porque los SSNN sin determinante, ya sean nombres de materia, como
\forma{carne} en \exref{ex:08-carne}, o plurales de nombres contables, como
\forma{relatos} en \exref{ex:08-relatos}, no producen el efecto de
particularización necesario para delimitar al evento. Para que este efecto se
produzca, el SN que funciona como COD debe estar acotado, ya sea mediante un
determinante, como vimos en \exref{ex:08-gm2}, o mediante un cuantificador,
como se muestra en \exref{ex:08-varias}.

\begin{exe}
  \ex\label{ex:08-varias} García Márquez \textit{escribió} [varias novelas] a lo
      largo de su vida.
\end{exe}

Por otra parte, los verbos de movimiento pueden ser delimitados por un
complemento preposicional que indique el límite de su realización.

\begin{exe}
  \ex Forrest Gump \textit{corrió} [de costa a costa del país / hasta quedar
      agotado].
  \ex \textit{Caminó} [hasta la otra cuadra].
\end{exe}

\begin{ejercicio}\label{ej:08-delimitar}
Transforme las siguientes actividades en realizaciones añadiéndoles un
complemento que las delimite y luego forme oraciones con ellas, según se muestra
en el ejemplo~1.

\begin{exe}
  \ex \textit{Leer}
  \ex \textit{Estudiar}
  \ex \textit{Diseñar}
  \ex \textit{Ir}
  \ex \textit{Construir}
  \ex \textit{Marchar}
\end{exe}

\begin{resuelto}[Respuesta posible]
1. Leyó [una novela de César Aira].
\end{resuelto}
\end{ejercicio}

\subsection{Incidencia del sujeto}\label{sec:08-sujeto}

Así como los complementos del V (por ejemplo, el COD) determinan la clase
aspectual del evento, también el carácter del SN que funciona como sujeto puede
incidir en su delimitación. Veamos cómo.

\begin{exe}
  \ex\label{ex:08-lluvia} {}[La lluvia] \textit{cayó / caía} sobre la ciudad
      durante toda la tarde.
  \ex\label{ex:08-rayo} {}[Un rayo] \textit{cayó / *caía} sobre el auto y lo
      destruyó.
  \ex\label{ex:08-rayos} {}[Rayos] \textit{cayeron / caían} sobre la ciudad
      durante toda la tarde.
\end{exe}

En \exref{ex:08-lluvia}, con un sujeto que refiere a un nombre de materia, como
\forma{la lluvia}, y que, por lo tanto, es incontable,\footnote{Los sustantivos
contables o discontinuos se refieren a objetos discretos y aislables, que no
pueden dividirse sin dejar de ser lo que son ---\forma{mesa, libro, coche}---.
Los sustantivos incontables o continuos tienen una constitución homogénea, por
lo que pueden dividirse sin perder su naturaleza ---\forma{agua, café, arena}---
\parencite[30]{giammatteo2009}.} el verbo \forma{caer} puede interpretarse como
un evento delimitado cuando está en perfecto simple ---\forma{cayó}--- y como no
delimitado y durativo cuando está en imperfecto ---\forma{caía}---. En cambio,
en \exref{ex:08-rayo}, con el sujeto \forma{un rayo}, que es un nombre contable,
el verbo solo admite perfecto simple y la única interpretación posible es la de
que se trata de un logro. Pero ¿qué sucede en \exref{ex:08-rayos}? ¿Por qué
ahora \forma{rayos} es compatible con el pretérito imperfecto? Según acabamos de
ver en relación con el COD, cuando los SSNN en plural no están determinados, el
evento no queda delimitado y admite interpretación durativa, lo que hace a
\forma{rayos} compatible con el pretérito imperfecto (v.~§\ref{sec:08-perfectivos}).

No solo la pluralidad morfológica del sujeto incide en la delimitación del
evento, ya que los nombres colectivos, que manifiestan pluralidad
semántica,\footnote{Los sustantivos individuales, en singular, se refieren a un
solo objeto y, si mencionan más de uno, deben pluralizarse ---\forma{tres libros,
muchos faroles, algunos ajíes}---. Los sustantivos colectivos, en forma
singular, designan varias entidades ---\forma{enjambre, rebaño, docena}---
\parencite[31]{giammatteo2009}.} también producen un efecto similar, de modo que
permiten una interpretación durativa de los eventos \exref{ex:08-multitud}, que
los individuales no admiten \exref{ex:08-profesor}.\pendiente{El original remite
dos veces a (16)}

\begin{exe}
  \ex\label{ex:08-multitud} {}[La multitud] \textit{salió} del estadio
      ordenadamente durante una hora.
  \ex\label{ex:08-profesor} {}[El profesor] \textit{salió} de su casa (*durante
      una hora).
\end{exe}

Asimismo, también el papel temático del sujeto puede alterar la interpretación
eventiva del predicado. Por ejemplo, \forma{rodear} en \exref{ex:08-parque} y
\forma{envolver} en \exref{ex:08-atmosfera}, con sujetos tema, se entienden como
verbos de estado, mientras que en \exref{ex:08-policia} y
\exref{ex:08-vendedor}, con sujeto agente, se interpretan como verbos de acción
del tipo eventivo de realización.

\begin{exe}
  \ex\label{ex:08-parque} {}[El parque] \textit{rodea} la casa.
  \ex\label{ex:08-policia} {}[La policía] \textit{rodeó} la manzana en 20
      segundos.
  \ex\label{ex:08-atmosfera} {}[La atmósfera] \textit{envuelve} a la Tierra.
  \ex\label{ex:08-vendedor} Con mucho cuidado [el vendedor] \textit{envolvió} el
      regalo para la Navidad.
\end{exe}

\begin{ejercicio}\label{ej:08-sujeto}
Compare las siguientes oraciones y explique los contrastes aspectuales que se
dan en cada grupo, según se muestra en el ejemplo resuelto~1.

\begin{exe}
  \ex
    \begin{xlist}
      \ex La luz del sol \textit{rebotó / rebotaba} en la pared toda la tarde.
      \ex La pelota \textit{rebotó / rebotaba} en la pared toda la tarde.
      \ex La bala \textit{rebotó} en la pared (*toda la tarde).
    \end{xlist}
  \ex
    \begin{xlist}
      \ex El soldado \textit{murió} durante la emboscada.
      \ex Juan \textit{murió} (*durante toda la tarde).
    \end{xlist}
  \ex
    \begin{xlist}
      \ex El agua líquida \textit{cubre} aproximadamente tres cuartas partes de
          la superficie terrestre.
      \ex Pedro \textit{cubrió} el auto para que el granizo no lo estropee.
    \end{xlist}
\end{exe}

\begin{resuelto}
1a. Con un sujeto que refiere a un nombre de materia incontable ---\forma{la luz
del sol}---, el verbo \forma{rebotar} en perfecto simple ---\forma{rebotó}--- se
interpreta como un evento delimitado, pero en imperfecto ---\forma{rebotaba}---
se considera no delimitado y durativo. 1b. Con el sujeto \forma{la pelota}, que
es un nombre común y contable, el verbo admite perfecto simple, con
interpretación de evento delimitado y puntual (= \sig{dio un rebote}); pero en
imperfecto el evento se considera reiterativo. 1c. Con el sujeto \forma{la
bala}, \forma{rebotar} se vuelve un evento semelfactivo, único e irrepetible,
por lo que rechaza el modificador durativo \forma{toda la tarde}.
\end{resuelto}
\end{ejercicio}

\subsection{Compatibilidad con distintas clases de modificadores temporales (MT)}\label{sec:08-mt}

Si comparamos los modificadores temporales (MMTT) de las siguientes oraciones,
vemos que en \exref{ex:08-estreno} y \exref{ex:08-regreso} refieren a un momento
de tiempo determinado respecto del que el evento se sitúa ---\forma{la semana
pasada, en 1994, a las nueve de la noche}---. En \exref{ex:08-duro} y
\exref{ex:08-tarea}, en cambio, el tiempo es concebido en su transcurrir
---\forma{dos horas, toda la tarde}---, sin atender a su culminación; así, por
ejemplo, \exref{ex:08-tarea-a} acepta continuarse con \forma{\ldots{}pero no la
terminamos}, como se muestra en \exref{ex:08-tarea-b}. Los modificadores
temporales de \exref{ex:08-estreno} y \exref{ex:08-regreso} se denominan
\termino{puntuales} (MTP), mientras que los de \exref{ex:08-duro} y
\exref{ex:08-tarea} son modificadores \termino{de duración} (MTD).

\begin{exe}
  \ex\label{ex:08-estreno} La película \textit{se estrenó} [la semana pasada / en
      1994].
  \ex\label{ex:08-regreso} \textit{Regresó} del cine [a las nueve de la noche].
  \ex\label{ex:08-duro} La proyección \textit{duró} [dos horas].
  \ex\label{ex:08-tarea}
    \begin{xlist}
      \ex\label{ex:08-tarea-a} \textit{Estuvimos haciendo} la tarea [toda la
          tarde].
      \ex\label{ex:08-tarea-b} \textit{Estuvimos haciendo} la tarea [toda la
          tarde], pero no la terminamos.
    \end{xlist}
\end{exe}

Por su parte, los modificadores que aparecen en \exref{ex:08-guion} y
\exref{ex:08-manejar}, si bien implican duración, delimitan un período de tiempo
y presentan al suceso como concluido. Se conocen como modificadores temporales
\termino{de período} (MTPer) y aceptan una paráfrasis o equivalencia con
\forma{tardar} + SN cuantificado + \forma{en} + infinitivo
\exref{ex:08-guion-b} o con \forma{llevarle (a alguien)} + SN cuantificado +
infinitivo \exref{ex:08-guion-c}.

\begin{exe}
  \ex\label{ex:08-guion}
    \begin{xlist}
      \ex \textit{Terminó} el guion de la película [en seis meses].
      \ex\label{ex:08-guion-b} \textit{Tardó} seis meses en terminar el trabajo.
      \ex\label{ex:08-guion-c} Le \textit{llevó} seis meses escribir el guion de
          la película.
    \end{xlist}
  \ex\label{ex:08-manejar} Aprendió a manejar [en un mes].
\end{exe}

Según veremos, las distintas subclases eventivas de verbos muestran afinidades e
incompatibilidades con los tres tipos de modificadores temporales mencionados.
Sin embargo, como sostiene \textcite[3047]{demiguel1999}, «al tratar de
información aspectual no se puede hablar de incompatibilidades o restricciones
auténticas, sino más bien de preferencias o tendencias», ya que, según veremos
en §\ref{sec:08-coercion}, las distintas subclases pueden cambiar su tipo
eventivo para compatibilizar con los modificadores u otros constituyentes
oracionales.

Las actividades [+din, +dur, $-$tél] y las realizaciones [+din, +dur, +tél]
presentan un comportamiento complementario respecto de los MTD y los MTPer: las
primeras típicamente se combinan con MTD \exref{ex:08-forrest}, mientras que las
segundas preferentemente seleccionan MTPer \exref{ex:08-recaudo}:

\begin{exe}
  \ex\label{ex:08-forrest} Forrest Gump \textit{corrió} de costa a costa (durante
      tres años / por mucho tiempo / *en una semana).
  \ex\label{ex:08-recaudo} La película \textit{recaudó} millones de dólares (en
      una semana / ?durante tres años / *por mucho tiempo).
\end{exe}

En cuanto a los estados [$-$din, +dur, $-$tél], una de sus características es su
duración potencialmente ilimitada, que impide que, por exceso, puedan emplearse
con MTD que los acoten \exref{ex:08-camarones}. Sin embargo, los estados
potencialmente finitos o transitorios son compatibles con MTD
\exref{ex:08-dan}:

\begin{exe}
  \ex[*]{\label{ex:08-camarones} Los camarones \textit{son} crustáceos [durante
      este año].}
  \ex\label{ex:08-dan} El teniente Dan \textit{era} el líder del pelotón de
      Forrest Gump [durante la guerra de Vietnam].
\end{exe}

\begin{recuadro}[Para ampliar]
En la bibliografía se suele explicar la diferencia entre estados de duración
ilimitada y estados finitos mediante la distinción que introduce
\textcite{carlson1977} entre \textit{individual-level predicates} (predicados de
nivel individual) y \textit{stage-level predicates} (predicados de estadio). Los
primeros son aquellos que persisten en el tiempo y pueden caracterizar una
entidad, como \forma{ser alto, tener ojos verdes} o \forma{ser cordobés}. Los
segundos son aquellos que pueden cambiar con el tiempo, como \forma{estar
enfermo, tener fiebre} o \forma{estar contento}. En español, esta distinción
entre \termino{estados de nivel individual} y \termino{estados de estadio}, en
líneas generales, se puede correlacionar con el uso de los verbos \forma{ser},
para los primeros ---\forma{Guillermo es enfermo desde su juventud}---, y
\forma{estar}, para los segundos ---\forma{Guillermo estuvo enfermo la semana
pasada}---.\footnote{Si bien la distinción entre \forma{ser} y \forma{estar}
tiene una base aspectual, existen varias excepciones en las que no podemos
extendernos aquí. Para ampliar, remitimos a \textcite[313 y
ss.]{bosque2009}.}
\end{recuadro}

Los logros [+din, $-$dur, +tél],\pendiente{El original dice aquí
«[+din, --dur, --tél]»} por su carácter no durativo, aceptan MTP referidos a un
momento muy breve \exref{ex:08-valla} o a un tiempo determinado
\exref{ex:08-1994}, mientras que los estados los rechazan
\exref{ex:08-ocho}.\pendiente{El original habla aquí de «MTPer»}

\begin{exe}
  \ex\label{ex:08-valla} Los refugiados \textit{saltaron} la valla [en un
      segundo].
  \ex\label{ex:08-1994} La película \textit{se estrenó} [en 1994].
  \ex[*]{\label{ex:08-ocho} Los camarones \textit{viven} en el mar [a las ocho].}
\end{exe}

Los semelfactivos típicamente se combinan con modificadores de repetición
(MRep) \exref{ex:08-tosio} y de frecuencia (MFr) \exref{ex:08-bostezaba}, pero
este tipo de modificadores son rechazados por los predicados de estado que, al
no tener una duración acotada, no admiten repetición ---\exref{ex:08-portero} y
\exref{ex:08-enferma-a}---, excepto que se trate de estados transitorios
\exref{ex:08-enferma-b}:

\begin{exe}
  \ex\label{ex:08-tosio} \textit{Tosió} [dos veces seguidas].
  \ex\label{ex:08-bostezaba} \textit{Bostezaba} [cada diez segundos].
  \ex[*]{\label{ex:08-portero} El portero de la escuela \textit{pertenece /
      perteneció} a una familia muy humilde [habitualmente / tres veces / en
      varias ocasiones].}
  \ex
    \begin{xlist}
      \ex[*]{\label{ex:08-enferma-a} [Habitualmente / Cada dos por tres] mi
          hermana \textit{es enferma} en invierno.}
      \ex\label{ex:08-enferma-b} {}[Habitualmente / Cada dos por tres] mi hermana
          \textit{está enferma} en invierno.
    \end{xlist}
\end{exe}

\subsection{La reinterpretación de la clase eventiva}\label{sec:08-coercion}

Cuando un verbo de una subclase eventiva se combina con un modificador que exige
una reinterpretación de su significado, se produce una «acomodación», para la
que \textcite{pustejovsky1995} utiliza la denominación de
\termino{coerción} o cambio del tipo eventivo.

\begin{destacado}
\textcite[330]{bosque2009} definen la \termino{coerción} o \termino{coacción}
como la «[o]peración léxica que cambia el tipo aspectual (una o más propiedades
aspectuales) de una pieza léxica. Sin este cambio la combinación resultante
sería no convergente».
\end{destacado}

Algunos de estos casos de coerción se producen:

\begin{itemize}
  \item cuando las realizaciones, que son télicas [+tél], se usan con MTD, por
        lo que estos modificadores pasan a referirse al desarrollo del evento
        \exref{ex:08-escribio}, mientras que si las actividades [$-$tél] se
        combinan con MTPer, estos se acomodan a hacer referencia al proceso o
        período de realización previo a la culminación del evento
        \exref{ex:08-vistio};
\end{itemize}

\begin{exe}
  \ex\label{ex:08-escribio} \textit{Escribió} el guion [durante un año].
  \ex\label{ex:08-vistio} \textit{Se vistió} [en quince minutos] y salió a la
      calle.
\end{exe}

\begin{itemize}
  \item cuando los logros se combinan con MTD, estos pueden hacer referencia al
        proceso previo \exref{ex:08-hirvio} o al estado resultante del evento
        \exref{ex:08-sefue}, o bien pueden tomar valor habitual
        \exref{ex:08-fuegos};
\end{itemize}

\begin{exe}
  \ex\label{ex:08-hirvio} La leche \textit{hirvió} [durante media hora].
  \ex\label{ex:08-sefue} \textit{Se fue} [por diez minutos].
  \ex\label{ex:08-fuegos} [Durante la semana de festejos] \textit{se lanzaron}
      fuegos artificiales.
\end{exe}

\begin{itemize}
  \item cuando logros y semelfactivos se combinan con MTPer o con MTD adquieren
        interpretación iterativa o habitual.
\end{itemize}

\begin{exe}
  \ex En China, Forrest \textit{ganó} varios partidos de pimpón [en una semana].
  \ex Un espectador de la tercera fila \textit{bostezó} sin parar [durante toda
      la película].
\end{exe}

\begin{ejercicio}\label{ej:08-construir}
Construya oraciones según lo pedido y luego indique la interpretación
resultante, según se muestra en los ejemplos resueltos de abajo.

\begin{exe}
  \ex Un verbo de estado con un MTD.
  \ex Un verbo de logro con un MTD.
  \ex Un verbo de realización con un MTPer.
  \ex Un verbo de estado transitorio con un MRep.
  \ex Un verbo de actividad con un MTD.
  \ex Un verbo semelfactivo con un MRep.
  \ex Un verbo de realización con un MTD.
  \ex Un verbo de logro con un MTPer.
\end{exe}

\begin{resuelto}[Respuestas posibles]
1. Los dinosaurios \textit{han existido} [por muchos millones de años]. El MTD
refuerza el valor durativo del estado. 2. \textit{Encontraron} muchos fósiles
[durante la excavación]. El MTD fuerza una interpretación durativa de repetición
(= \sig{hubo muchos hallazgos}) del logro.
\end{resuelto}
\end{ejercicio}

\subsection{Compatibilidad de los estados con otras clases de modificadores verbales}\label{sec:08-estados}

Dado su carácter homogéneo y que no indican cambio interno en su desarrollo, los
estados no son compatibles con modificadores de dinamicidad, como los de
manera, del tipo de \forma{rápidamente, con alegría, sin miedo}, ni con los
modificadores aspectuales que indican grado de completamiento del evento, como
\forma{poco a poco, de a poco, gradualmente, hasta la mitad}.

\begin{exe}
  \ex[*]{Las ballenas \textit{son} cetáceos [rápidamente / sin miedo].}
  \ex[*]{Indiana Jones \textit{teme} a las víboras [poco a poco / hasta la
      mitad].}
\end{exe}

Como los estados carecen de sujeto agente, tampoco son compatibles con los
adverbios orientados al sujeto, que describen el modo o manera en que el agente
realiza la acción ---\forma{cuidadosamente, esforzadamente,
inconscientemente, inteligentemente}, etc.---.

\begin{exe}
  \ex[*]{Los triángulos \textit{tienen} tres lados y tres ángulos
      [inteligentemente / deliberadamente].}
  \ex[*]{Algunos alumnos \textit{saben} inglés [cuidadosamente].}
\end{exe}

\subsection{Interacción del tipo eventivo (aspecto léxico) con el aspecto y el modo gramatical}\label{sec:08-interaccion}

\subsubsection{Tiempos perfectivos e imperfectivos}\label{sec:08-perfectivos}

Los \termino{tiempos imperfectivos} resaltan el carácter durativo y no
conclusivo de los eventos no télicos, como los estados \exref{ex:08-josefa} y
las actividades \exref{ex:08-mario}.

\begin{exe}
  \ex\label{ex:08-josefa} Josefa \textit{vivía} en Nueva York durante su
      juventud.
  \ex\label{ex:08-mario} Mario \textit{trabajaba} en una carpintería de
      Avellaneda.
\end{exe}

En cuanto a los eventos télicos ---realizaciones y logros---, la llamada
\termino{paradoja imperfectiva} hace referencia a que, cuando están en un tiempo
imperfectivo, no implican una culminación, \exref{ex:08-leia} y
\exref{ex:08-llegaba}, a diferencia de lo que sucede con los estados y
actividades, \exref{ex:08-vivian} y \exref{ex:08-corria}, que implican que el
suceso efectivamente ocurrió.

\begin{exe}
  \ex\label{ex:08-leia} Juan \textit{leía} una novela de misterio $\not\Rightarrow$
      Juan \textit{leyó} una novela de misterio.
  \ex\label{ex:08-llegaba} El doctor López \textit{llegaba} a su casa
      $\not\Rightarrow$ El doctor López \textit{llegó} a su casa.
  \ex\label{ex:08-vivian} Mis amigos \textit{vivían} en Hurlingham $\Rightarrow$
      Mis amigos \textit{vivieron} en Hurlingham.
  \ex\label{ex:08-corria} Forrest \textit{corría} por las rutas $\Rightarrow$
      Forrest \textit{corrió} por las rutas.
\end{exe}

Los tiempos imperfectivos son los preferidos por cualquier clase eventiva para
expresar hechos habituales o repetitivos.

\begin{exe}
  \ex \textit{Viajaba} todos los fines de semana al interior del país.
  \ex Estaba tan nervioso que \textit{entraba} y \textit{salía} continuamente.
\end{exe}

Por su parte, los \termino{tiempos perfectivos}, como el perfecto simple, se
emplean con cualquier clase de evento para poner de relieve su terminación.

\begin{exe}
  \ex Mi jefe \textit{leyó} atentamente el informe. (realización)
  \ex Marina \textit{entró} cansada a la oficina. (logro)
  \ex Forrest \textit{sufrió} mucho en su infancia. (estado)
  \ex El sol \textit{brilló} durante toda la semana. (actividad)
\end{exe}

También prefieren los tiempos perfectivos los eventos únicos:

\begin{exe}
  \ex Noam Chomsky \textit{nació} el 7 de diciembre de 1928.
\end{exe}

\subsubsection{La perífrasis progresiva}\label{sec:08-progresiva}

Dado que las perífrasis progresivas\footnote{Las perífrasis verbales se tratan
en el \capref{11}.} indican «el proceso de cambio asociado a una eventualidad»,
resultan incompatibles con los estados, que se caracterizan por no cambiar o no
avanzar en el tiempo \parencite[301]{bosque2009} ---\exref{ex:08-triangulo} y
\exref{ex:08-tigres}---. Sin embargo, los estados potencialmente delimitables
permiten la inclusión de un participante que se esfuerza en el sostenimiento de
la situación, que se presenta como más dinámica y compatible con un tiempo
expresado por una perífrasis progresiva ---\exref{ex:08-blanda} y
\exref{ex:08-gustando}---:

\begin{exe}
  \ex[*]{\label{ex:08-triangulo} Este triángulo \textit{está teniendo} tres
      lados.}
  \ex[*]{\label{ex:08-tigres} Los tigres \textit{están siendo} felinos.}
  \ex\label{ex:08-blanda} Marina \textit{está siendo} muy blanda con su hija.
  \ex\label{ex:08-gustando} Me \textit{está gustando} mucho la novela que me
      prestaste.
\end{exe}

Los logros tampoco aceptan los tiempos progresivos porque, al ser instantáneos,
no pueden indicar un desarrollo temporal del evento. Sin embargo, como ya hemos
señalado, admiten una referencia tanto al tiempo previo \exref{ex:08-resolviendo}
como al posterior \exref{ex:08-hirviendo}:

\begin{exe}
  \ex\label{ex:08-resolviendo} \textit{Estuvo resolviendo} el ejercicio un buen
      rato. (= \sig{Estuvo haciéndolo un buen rato hasta que lo resolvió})
  \ex\label{ex:08-hirviendo} La mezcla \textit{estuvo hirviendo} media hora más,
      como lo indicaba la receta. (= \sig{Continuó hirviendo media hora después
      del momento de ebullición})
\end{exe}

\subsubsection{Perífrasis fasales}\label{sec:08-fasales}

Las perífrasis fasales focalizan distintos momentos del evento: ya sea el inicio
---\forma{empezar / comenzar a}---, el fin ---\forma{dejar / acabar / terminar
de}--- o el desarrollo ---\forma{continuar / seguir}--- y, para hacerlo, exigen
que el verbo principal sea durativo y tenga delimitación. En consecuencia, los
verbos de logro no pueden emplearse con estas perífrasis \exref{ex:08-llegar},
excepto que tengan interpretación habitual \exref{ex:08-coleccionar}. Por la
misma razón interpretativa, los semelfactivos, que indican repetición del
evento, son posibles con estas perífrasis \exref{ex:08-bostezar}.

\begin{exe}
  \ex[*]{\label{ex:08-llegar} \textit{Empezó a / terminó de llegar / continuó
      llegando} toda la tarde.}
  \ex\label{ex:08-coleccionar} \textit{Empezó a / dejó de coleccionar / continuó
      coleccionando} objetos extraños.
  \ex\label{ex:08-bostezar} \textit{Empezó a / dejó de bostezar} enseguida /
      \textit{continuó bostezando} durante toda la función.
\end{exe}

Por otra parte, dado que los estados suelen carecer de delimitación, no pueden
utilizarse ni con las perífrasis inceptivas ---\forma{empezar / comenzar a}---
\exref{ex:08-salon} ni con las culminativas ---\forma{terminar de, acabar de}---
\exref{ex:08-contenta-a}, aunque admiten la perífrasis cesativa con \forma{dejar
de} \exref{ex:08-contenta-b}. La diferencia, para \textcite[95]{cinque1999},
estriba en que el \termino{aspecto cesativo} caracteriza una situación que ha
quedado inconclusa. Así, por ejemplo, en \forma{dejó de lavar los platos}, el
evento se interrumpe sin que necesariamente se haya completado, es decir, sin
que se hayan terminado de lavar todos los platos. En cambio, en el
\termino{aspecto culminativo}, el evento sí alcanza su punto final natural. Así,
en \forma{terminó de lavar los platos}, se interpreta que el evento culmina
cuando no quedan más platos por lavar.

\begin{exe}
  \ex[*]{\label{ex:08-salon} El salón \textit{empezó / comenzó a medir} 20
      metros.}
  \ex
    \begin{xlist}
      \ex[*]{\label{ex:08-contenta-a} \textit{Terminó de estar contenta} cuando
          supo la verdad.}
      \ex\label{ex:08-contenta-b} \textit{Dejó de estar contenta} cuando supo la
          verdad.
    \end{xlist}
\end{exe}

El hecho de que los estados no presenten fases de desenvolvimiento no favorece
la utilización de estos predicados con las perífrasis inceptivas; sin embargo,
no lo excluye totalmente \exref{ex:08-corpulento}. Además, los estados
transitorios admiten esta combinación sin problemas \exref{ex:08-recuperar}.

\begin{exe}
  \ex\label{ex:08-corpulento} \textit{Empezó a ser corpulento} a los doce años,
      cuando se desarrolló.
  \ex\label{ex:08-recuperar} \textit{Empezó a estar contenta} cuando supo que se
      iba a recuperar de su enfermedad.
\end{exe}

\subsubsection{Modo imperativo}\label{sec:08-imperativo}

La no agentividad de los estados, que deriva de su carácter estable y homogéneo
(sin fases), tiene como consecuencia que no se empleen en imperativo,
\exref{ex:08-ojos} y \exref{ex:08-alto}, excepto cuando representan un
comportamiento que puede concebirse como dependiente de la voluntad,
\exref{ex:08-caritativo} y \exref{ex:08-quieto}.

\begin{exe}
  \ex[*]{\label{ex:08-ojos} \textit{Tené} ojos claros.}
  \ex[*]{\label{ex:08-alto} \textit{Sé} alto.}
  \ex\label{ex:08-caritativo} \textit{Sé} caritativo y ayudá a quien lo necesita.
  \ex\label{ex:08-quieto} \textit{Quedate} quieto.
\end{exe}

\begin{ejercicio}\label{ej:08-naufrago}
Lea atentamente el siguiente fragmento adaptado de \textit{Relato de un
náufrago}, de G.~García Márquez (1970), y luego complete la tabla de abajo con
lo pedido, según se muestra en los ejemplos ya resueltos.

% Fragmento literario protegido: se incorpora a mano desde el original.
\textoaincorporar{el fragmento adaptado de la presentación de \textit{Relato de
un náufrago}, desde «El 28 de febrero de 1955 se conoció la noticia» hasta
«\ldots{}sus diez días en el mar», con los predicados que figuran en la tabla}

\begin{consignas}
  \item Determine la clase eventiva de todos los predicados del fragmento.
  \item Identifique, si corresponde, los modificadores temporales que reciben.
  \item Reconozca el aspecto gramatical de cada verbo y determine su incidencia
        en cada predicación.
  \item En todos los casos, especifique si la lectura o interpretación
        resultante es puntual o durativa.
  \item Determine si hay compatibilidad o coerción (cambio de tipo eventivo)
        entre los distintos factores considerados.
\end{consignas}

{\footnotesize
\begin{tabularx}{\linewidth}{@{}>{\raggedright\arraybackslash}p{19mm}
    >{\raggedright\arraybackslash}p{13mm} L L
    >{\raggedright\arraybackslash}p{10mm} L@{}}
  \toprule
  Predicado & Clase eventiva & MT & Aspecto verbal & Lectura &
    Compatibilidad / coerción \\
  \midrule
  \forma{se conoció} & estado & MTP \forma{el 28 de febrero de 1955} &
    perfectivo: presenta al evento como concluido & puntual & coerción del tipo
    eventivo por el MT y el aspecto verbal \\
  \forma{habían caído y desaparecido} & logro & --- & perfectivo: refuerza el
    valor conclusivo & puntual & compatibilidad entre tipo eventivo y aspecto
    verbal \\
  \forma{viajaba} & & & & & \\
  \forma{había sido sometida a reparaciones} & & & & & \\
  \forma{llegó} & & & & & \\
  \forma{se inició} & & & & & \\
  \forma{hacen oficios} & & & & & \\
  \forma{se desistió} & & & & & \\
  \forma{fueron declarados} & & & & & \\
  \forma{apareció} & & & & & \\
  \forma{había permanecido} & & & & & \\
  \forma{se llamaba} & & & & & \\
  \forma{es} & & & & & \\
  \forma{contó} & & & & & \\
  \forma{fue publicada} & & & & & \\
  \forma{tomaba notas} y \forma{soltaba preguntas} & & & & & \\
  \forma{pudimos reconstruir el relato} & & & & & \\
  \bottomrule
\end{tabularx}}
\end{ejercicio}

\begin{recuadro}[Reflexión final]
La estructura eventiva, dependiente tanto del significado léxico del verbo como
de la categoría de aspecto, se muestra como un fenómeno esencialmente
combinatorio y con múltiples aristas. Es así que, para todo predicado, hay que
tener en cuenta no solo la naturaleza del verbo, es decir, la clase aspectual a
la que pertenece (estado, actividad, realización, logro o semelfactivo), sino
también la estructura sintáctica que conforma con los distintos modificadores
que recibe, incluido el sujeto, y los distintos tiempos y formas verbales en que
se manifiesta el verbo. Si bien se trata de una cuestión compleja, su
reconocimiento resulta crucial para identificar, a partir de la naturaleza del
verbo y la configuración morfosintáctica en que se manifiesta, cuál será la
lectura o interpretación semántica que corresponde al predicado.
\end{recuadro}

% --- Secciones finales -------------------------------------------------
\lecturasrecomendadas

\solucionario

\begin{solucion}{\ref{ej:08-clases}}
\begin{center}\small
\begin{tabular}{@{}lcccl@{}}
  \toprule
  Verbo & Duratividad & Dinamicidad & Telicidad & Clase eventiva \\
  \midrule
  \forma{buscar}                         & + & + & $-$ & actividad \\
  \forma{dirigir una película}           & + & + & + & realización \\
  \forma{vivir una serie de aventuras}   & + & $-$ & +\textsuperscript{a} & actividad \\
  \forma{recuperar}                      & $-$ & + & + & logro \\
  \forma{enfrentar}                      & + & + & $-$ & actividad \\
  \forma{golpear}                        & $-$\textsuperscript{b} & + & $-$ & semelfactivo \\
  \forma{derribar}                       & $-$ & + & + & logro \\
  \forma{atrapar}                        & $-$ & + & + & logro \\
  \forma{tener}                          & + & $-$ & $-$ & estado \\
  \forma{adquirir}                       & $-$ & + & + & logro \\
  \forma{asegurar}                       & $-$ & + & + & logro \\
  \forma{escribir el libro}              & + & + & + & realización \\
  \bottomrule
\end{tabular}
\end{center}
\textsuperscript{a} El evento \forma{vivir una serie de aventuras} está
delimitado por el modificador temporal encabezado por \forma{hasta}.
\textsuperscript{b} Este evento se vuelve durativo por la repetición.
\pendiente{Los rasgos de \textit{vivir una serie de aventuras} ($-$din, +tél) no
corresponden a una actividad}
\end{solucion}

\begin{solucion}{\ref{ej:08-delimitar}}
Ejemplos de respuestas posibles. 2. Estudiamos [la voz pasiva] para rendir el
examen. 3. Diseñaron [un nuevo espacio verde] para la ciudad. 4. Fuimos [hasta
la plaza]. 5. ¿Ya construyeron [la nueva autopista]? 6. Marchó [hasta la
frontera].
\end{solucion}

\begin{solucion}{\ref{ej:08-sujeto}}
2a. El verbo de logro \forma{morir} es compatible con \forma{durante la
emboscada}, que limita la ocurrencia del evento a un período de tiempo. 2b.
\forma{Durante toda la tarde} indica duración y es incompatible con un verbo de
logro, puntual y delimitado, como \forma{morir}. 3. En 3a, con sujeto tema,
\forma{cubrir} se entiende como un verbo de estado, mientras que en 3b, con
sujeto agente, se interpreta como verbo de acción, del tipo eventivo de
realización.
\end{solucion}

\begin{solucion}{\ref{ej:08-construir}}
Ejemplos de respuestas posibles. 3. \textit{Planchó} toda la ropa [en una hora].
El MTPer mide la duración de la realización. 4. \textit{Estuvo} [varias veces] en
la ciudad. El MRep refiere a la reiteración de un estado transitorio. 5. Todas
las mañanas \textit{corre} [media hora] alrededor de la plaza. El MTD mide la
duración del evento de actividad. 6. \textit{Tosió} [tres o cuatro veces] durante
la conferencia. El MRep indica la repetición del evento semelfactivo. 7. Los
turistas \textit{recorrieron} la ciudad [durante la tarde]. El MTD mide el
tiempo del evento de realización. 8. \textit{Llegó} de su casa a la Facultad [en
dos horas]. El MTPer mide el tiempo previo al evento.
\end{solucion}

\begin{solucion}{\ref{ej:08-naufrago}}
{\footnotesize\setlength{\tabcolsep}{3pt}
\begin{longtable}{@{}>{\raggedright\arraybackslash}p{18mm}
    >{\raggedright\arraybackslash}p{14mm} >{\raggedright\arraybackslash}p{20mm}
    >{\raggedright\arraybackslash}p{20mm} >{\raggedright\arraybackslash}p{11mm}
    >{\raggedright\arraybackslash}p{32mm}@{}}
  \toprule
  Predicado & Clase eventiva & MT & Aspecto verbal & Lectura &
    Compatibilidad / coerción \\
  \midrule
  \endhead
  \forma{viajaba} & actividad & --- & imperfectivo: refuerza el valor durativo &
    durativa & compatibilidad entre tipo eventivo y aspecto verbal \\
  \forma{había sido sometida a reparaciones} & realización & --- &
    perfectivo\textsuperscript{a} & durativa & compatibilidad entre tipo eventivo
    y aspecto verbal \\
  \forma{llegó} & logro & MTP \forma{dos horas después} & perfectivo: refuerza
    el valor conclusivo & puntual & compatibilidad entre tipo eventivo, MTP y
    aspecto verbal \\
  \forma{se inició} & logro & MTP \forma{de inmediato} & perfectivo: refuerza el
    valor conclusivo & puntual & compatibilidad entre tipo eventivo, MTP y
    aspecto verbal \\
  \forma{hacen oficios} & actividad & --- & imperfectivo: refuerza el valor
    durativo & habitual & compatibilidad entre tipo eventivo y aspecto verbal;
    incidencia del COD plural en la lectura habitual \\
  \forma{se desistió} & logro & MTP \forma{al cabo de cuatro días} & perfectivo:
    refuerza el valor conclusivo & puntual & compatibilidad entre tipo eventivo,
    MTP y aspecto verbal \\
  \forma{fueron declarados} & logro & --- & perfectivo: refuerza el valor
    conclusivo & puntual & compatibilidad entre tipo eventivo y aspecto verbal \\
  \forma{apareció} & logro & MTP \forma{una semana más tarde} & perfectivo:
    refuerza el valor conclusivo & puntual & compatibilidad entre tipo eventivo,
    MTP y aspecto verbal \\
  \forma{había permanecido} & estado & MTD \forma{diez días} & perfectivo:
    presenta al evento como concluido & durativa & la compatibilidad entre tipo
    eventivo y MTD produce la lectura durativa; el aspecto perfecto señala
    anterioridad y conclusividad del evento, pero no coerciona el tipo eventivo
    del verbo \\
  \forma{se llamaba} & estado & --- & imperfectivo: refuerza el valor durativo &
    durativa & compatibilidad entre tipo eventivo y aspecto verbal \\
  \forma{es} & estado & --- & imperfectivo: refuerza el valor durativo & durativa
    & compatibilidad entre tipo eventivo y aspecto verbal \\
  \forma{contó} & realización (delimitada por el COD \forma{lo}) & --- &
    perfectivo: refuerza el valor conclusivo & durativa & compatibilidad entre
    tipo eventivo y aspecto verbal \\
  \forma{fue publicada} & realización & MTP \forma{un mes después} & perfectivo:
    refuerza el valor conclusivo & puntual & compatibilidad entre tipo eventivo,
    MTP y aspecto verbal \\
  \forma{tomaba notas} y \forma{soltaba preguntas} & logro & MTD \forma{durante
    las cuales} (veinte sesiones de seis horas diarias) & imperfectivo: refuerza
    el valor durativo & durativa & coerción del tipo eventivo de logro por el MTD
    y el aspecto verbal, que producen la lectura durativa \\
  \forma{pudimos reconstruir el relato} & realización & MTPer \forma{en veinte
    sesiones de seis horas diarias} & perfectivo: refuerza el valor conclusivo &
    durativa & la compatibilidad entre tipo eventivo y MTPer produce la lectura
    durativa; el tiempo perfecto destaca la conclusión del evento, pero no
    coerciona el tipo eventivo del verbo \\
  \bottomrule
\end{longtable}}

\textsuperscript{a} El sustantivo plural \forma{reparaciones} no permite
concluir la delimitación del evento.
\end{solucion}
